\documentclass[12pt, a4paper]{article}

% Русский язык
\usepackage[T2A]{fontenc}
\usepackage[utf8]{inputenc}
\usepackage[english,russian]{babel}

% Математика
\usepackage{amsmath, amssymb}

% Таблицы
\usepackage{booktabs}
\usepackage{multirow}
\usepackage{array}
\usepackage{tabularx}

% Графика (если понадобится вставить рисунки)
\usepackage{graphicx}

% Гиперссылки и оглавление
\usepackage{hyperref}
\hypersetup{
    colorlinks=true,
    linkcolor=black,
    citecolor=black,
    urlcolor=black
}

% Поля (стандартные для ВКР/курсовых)
\usepackage[left=3cm, right=1.5cm, top=2cm, bottom=2cm]{geometry}

% Отступы и межстрочный интервал
\usepackage{setspace}
\onehalfspacing

% Нумерация страниц
\usepackage{lastpage}
\usepackage{fancyhdr}
\pagestyle{fancy}
\fancyhf{}
\cfoot{\thepage}

% Переопределение названий
\renewcommand{\contentsname}{Оглавление}

\begin{document}

% Титульный лист
\begin{titlepage}
    \centering
    {\scshape\Large НАЦИОНАЛЬНЫЙ ИССЛЕДОВАТЕЛЬСКИЙ УНИВЕРСИТЕТ \par}
    {\scshape\large «ВЫСШАЯ ШКОЛА ЭКОНОМИКИ» \par}
    \vspace{0.5cm}
    {\large Факультет экономических наук \par}
    \vspace{3cm}
    {\huge\bfseries Отношение потребителей к использованию искусственного интеллекта в дерматологии \par}
    \vspace{3cm}
    {\large Выполнили: \par}
    {\large студенты образовательной программы \par}
    {\large «Экономика и анализ данных» \par}
    {\large Дугаев Д.М., Космынин А.Э., Малакшанидзе А.Д. \par}
    \vfill
    {\large Научный руководитель: Егорова Л.Г. \par}
    \vspace{1cm}
    {\large Москва, 2026 \par}
\end{titlepage}

% Оглавление
\tableofcontents
\newpage

% Введение
\section{ВВЕДЕНИЕ}
\label{sec:intro}

Цифровая трансформация здравоохранения создает принципиально новый контекст для взаимодействия пациентов с медицинскими услугами. Искусственный интеллект (ИИ) стремительно проникает в клиническую практику: системы компьютерного зрения на основе сверточных нейронных сетей уже демонстрируют диагностическую точность до 97,6\% в задачах оценки тяжести акне \cite{Traini2025}, а мировой рынок ИИ в здравоохранении, оценивавшийся в 201 млрд долларов США в 2021 году, к 2030 году прогнозируется на уровне 1,2 трлн долларов \cite{Lamotkin2024}. Вместе с тем внедрение технологий наталкивается на серьезный барьер: готовность потребителей принять ИИ существенно расходится между странами и социальными группами \cite{Jutzi2020, McRae2025}.

Дерматология представляет собой особенно показательную область для изучения этой проблематики. С одной стороны, именно здесь ИИ достиг наиболее впечатляющих клинических результатов благодаря визуальной природе диагностики. С другой — дерматологические услуги сочетают в себе высокую субъективность эстетических ожиданий пациентов и выраженную социокультурную чувствительность, что делает анализ потребительских предпочтений особенно актуальным \cite{Artsrouni2025}.

Несмотря на растущий интерес к теме, существующие исследования обнаруживают принципиальные методологические ограничения. Большинство работ либо опираются на ретроспективный анализ технических характеристик алгоритмов \cite{Traini2025}, либо используют неэкономические опросные методы, не способные измерить готовность пациентов идти на компромисс между ценой, качеством и скоростью услуги \cite{Jutzi2020}. Применение метода дискретного выбора (DCE), позволяющего выявить реальные предпочтения через моделирование ситуаций выбора, остается редкостью в данной предметной области.

Настоящая работа направлена на восполнение этого пробела. Цель исследования — выявить структуру потребительских предпочтений в отношении форматов дерматологической консультации с участием ИИ и установить относительный вес ключевых атрибутов услуги при принятии решения о выборе.

Для достижения поставленной цели решаются следующие задачи:
\begin{enumerate}
    \item систематизировать теоретические подходы к анализу восприятия ИИ-технологий в медицине;
    \item разработать дизайн DCE-эксперимента, адекватный задаче измерения потребительских предпочтений в дерматологии;
    \item собрать первичные данные и провести их эконометрическую обработку;
    \item проверить совокупность гипотез о детерминантах выбора формата консультации;
    \item сформулировать содержательные выводы для исследователей и практиков.
\end{enumerate}

Объект исследования — российские потребители дерматологических услуг в возрасте от 18 до 27 лет. Предмет исследования — структура их предпочтений в отношении форматов консультации: «только живой врач», «только ИИ» и «врач совместно с ИИ» — в зависимости от атрибутов услуги (точности диагностики, времени ожидания, стоимости). При этом стоимость рассматривается как межгрупповой фактор рандомизации, а не как варьирующий атрибут внутри одного выбора респондента.

Методологическую основу работы составляют теория случайной полезности (Random Utility Theory) и метод дискретного выбора (DCE). Эмпирической базой служат данные онлайн-опроса 296 респондентов, собранные в марте 2026 года. Основным инструментом анализа является мультиномиальная логит-модель (MNL) с оценкой относительной важности атрибутов и взаимодействующих эффектов.

Работа состоит из введения, трех глав, заключения и списка использованных источников. В первой главе представлены теоретические основы и обзор литературы. Вторая глава описывает дизайн исследования, данные и методологию. Третья глава содержит эмпирические результаты и их обсуждение. Заключение подводит итоги работы, резюмируя основные результаты, выводы и ответы на поставленные исследовательские вопросы.

\newpage

% Глава 1
\section{ТЕОРЕТИЧЕСКИЕ ОСНОВЫ И ОБЗОР ЛИТЕРАТУРЫ}

\subsection{Экономические предпосылки и цифровизация здравоохранения}
Современные системы здравоохранения функционируют в условиях нарастающего экономического давления, обусловленного одновременным действием нескольких структурных факторов.

Во-первых, наблюдается глобальный кадровый дефицит: по нормативам ВОЗ для обеспечения базовой медицинской помощи требуется минимум 2,5 врача на 1 000 человек, однако в 45\% стран этот показатель составляет менее 1 на 1 000 \cite{Liu2021}.

Во-вторых, существует проблема с увеличением расходов на лечение тяжелых заболеваний, в частности в области онкодерматологии. За 2020 год зафиксировано более 1,5 млн новых случаев рака кожи, а к 2040 году прогнозируется рост заболеваемости меланомой до 500 000 новых случаев в год, что создает колоссальную нагрузку на бюджеты \cite{Jagemann2024}. Также огромные издержки вызваны средним уровнем диагностических ошибок: около 5,08\%, что затрагивает более 12 млн пациентов в США ежегодно \cite{Liu2021}.

В этих условиях искусственный интеллект (ИИ) рассматривается как механизм оптимизации аллокации ресурсов и снижения издержек. Объем мирового рынка ИИ в здравоохранении оценивался в 201 млрд долларов США уже в 2021 году, а к 2030 году прогнозируется рост до 1,2 трлн долларов \cite{Lamotkin2024}. Внедрение ИИ-алгоритмов (например, сверточных нейронных сетей, доказанная диагностическая точность которых достигает 97,6\% \cite{Traini2025}) решает проблему кадрового дефицита. Также оно ускоряет переход к медицине, ориентированной на индивидуальный подход и активное участие пациента — 4П-медицине (предиктивной, превентивной, персонализированной и партиципаторной) \cite{Myzrova2024}.

Однако для успешной интеграции ИИ необходимо учитывать не только технические метрики, но и готовность потребителей принять технологии, что требует глубокого экономического и социального анализа поведения общества.

\subsection{Теоретическая рамка: восприятие ИИ через призму теории полезности}
Традиционные социологические модели, такие как TAM \cite{Davis1989} и UTAUT \cite{Venkatesh2003}, не совсем точны для анализа восприятия обществом ИИ, так как они измеряют абстрактные предпочтения пользователей, не учитывая ограничения вроде времени и финансов, сопровождающих выбор медицинских услуг в реальности.

Данное исследование опирается на RUT (Random Utility Theory), согласно которой пациент выступает как рациональный агент, максимизирующий свою выгоду через оценку совокупности характеристик услуги. Решение об использовании ИИ формируется в компромиссе (trade-off) между клиническими выгодами (повышением точности), а также логистическими и экономическими издержками (временем ожидания, стоимостью).

Эта теория реализуется методологически через DCE (Discrete Choice Experiment): эксперимент предлагает реальные ситуации выбора, заставляющие респондента расставлять приоритеты (например, пожертвовать скоростью ради участия живого врача). Таким образом можно рассчитать относительную значимость каждого атрибута услуги при принятии решения \cite{Jagemann2024, Liu2021}.

\subsection{ИИ как ассистент, а не замена}
Эмпирические данные показывают значительную поддержку ИИ при условии его использования в гибридном формате — с врачом. При этом идея полной автоматизации решительно отвергается как медицинским сообществом, так и пациентами. Опрос 1 228 китайских дерматологов показал, что 95,36\% специалистов рассматривают ИИ исключительно для помощи в повседневной диагностике, и лишь 3,42\% допускают замену врача алгоритмом \cite{Shen2020}. Это согласуется с глобальным опросом (N = 1 271), где 77,3\% респондентов уверены в улучшении качества обслуживания с ИИ, однако страх потери работы зафиксирован лишь у 5,5\% \cite{Polesie2020}. В России из 301 опрошенного врача 85\% признают пользу ИИ: 79\% из них — для анализа данных, 74\% — для оптимизации процессов. При этом 76\% уверены, что специалисты, использующие ИИ, вытеснят тех, кто его игнорирует \cite{Orlova2023}.

С точки зрения пациентов: в немецком исследовании (N = 298) 94\% поддержали использование ИИ в качестве системы поддержки принятия решений, однако лишь 41\% согласились бы на диагностику исключительно машиной \cite{Jutzi2020}. По мониторинговым данным конца 2024 года в России 43\% населения верят, что ИИ улучшит качество лечения за счет снижения врачебных ошибок, и 53\% испытывали бы дискомфорт, если бы их лечащий врач делегировал постановку диагноза нейросети \cite{VCIOM2024}. Это расхождение подчеркивает необходимость изучения региональной и социокультурной специфики. Тем не менее для всех стран един фундаментальный консенсус: ИИ воспринимается исключительно как инструмент поддержки принятия решений. Глобальный опрос показывает, что в случае расхождения диагнозов врача и ИИ подавляющее большинство пациентов (84,4\%) предпочтут мнение живого специалиста \cite{Wu2024}. При этом интересно, что ИИ имеет доказанную клинико-психологическую ценность: исследование ИИ-трихоскопии показало, что визуализация данных алгоритма повышает доверие пациента врачу и мотивацию к лечению в 54,5\% случаев \cite{Kearney2025}.

\subsection{Дивергенция приоритетов и экономические факторы выбора (Trade-offs)}
Предпочтения пациентов формируются под влиянием стоимости, точности, времени и формата предоставления услуги. Исследования, использующие методологию DCE, наглядно демонстрируют весовые коэффициенты этих факторов. Точность диагностики является главным приоритетом: эксперимент в Китае (N = 428) выявил, что точность имеет наибольший удельный вес (39,29\%) при выборе метода диагностики \cite{Liu2021}. Экономический фактор занимает второе место по значимости (21,69\%) \cite{Liu2021}. Пациенты склонны предпочитать ИИ, если это снижает их личные расходы (out-of-pocket costs), однако готовы доплатить за гибридный формат. Немецкое исследование (N = 126) по скринингу рака кожи также подтвердило, что формат «ИИ + Врач» критически важен, но чувствительность к цене сохраняется \cite{Jagemann2024}. Время ожидания является сильным преимуществом ИИ: снижение времени ожидания приема увеличивает привлекательность ИИ-решений, особенно для рутинного скрининга (вес фактора 29,8\%) \cite{Jagemann2024, Liu2021}.

Вопросы автономии и ответственности порождают конфликт: пациенты настаивают на праве самостоятельно решать, будет ли ИИ вовлечен в их лечение, значительно сильнее врачей (4,2 против 3,3 балл из 5) \cite{Garcia2025}. Также 59,5\% российских врачей готовы брать юридическую ответственность на себя, однако 31,9\% считают, что финансово и юридически за ошибку алгоритма должен отвечать разработчик \cite{Orlova2023}.

\subsection{Этические и алгоритмические ограничения: барьеры для доверия}
Доверие потребителей напрямую зависит от прозрачности и объективности технологии. Барьером выступает проблема «черного ящика» и смещения данных (bias). Систематический анализ показал, что крупнейшие открытые датасеты (например, ACNE04) преимущественно состоят из фенотипов восточноазиатских пациентов \cite{Traini2025}. Нерепрезентативность выборок по шкале фототипов кожи ведет к этнолукизму — дискриминации и снижению качества диагностики из-за этнических особенностей внешности и фототипа кожи. Это усугубляет проблему экономического и социального неравенства в доступе к качественной помощи \cite{Artsrouni2025}.

Из-за повышенного риска в виде предвзятости данных (bias) и эффекта «черного ящика» готовность платить (WTP) рационального потребителя падает: он либо требует существенного снижения цены, либо вовсе отказывается от услуги. В качестве решения рассматривается объяснимый ИИ (Explainable AI), интерпретирующий свои выводы. Так, за счет прозрачности субъективные риски пациента снижаются, а его доверие и готовность платить растут \cite{Myzrova2024, Jagemann2024}.

\subsection{Пробелы в исследованиях}
Анализ литературы доказывает расположенность потребителей к гибридной медицине, однако выявляет существенные пробелы. Большинство существующих работ опираются либо на ретроспективный анализ изолированных датасетов \cite{Traini2025}, либо на неэкономические кросс-секционные опросы \cite{Jutzi2020, Shen2020}, не измеряющие готовность пациентов идти на компромисс между ценой, качеством и скоростью. Совсем немногие исследования применяют методологию DCE для оценки экономической эластичности спроса на услуги ИИ-дерматологии с учетом региональной специфики и уровня цифровой грамотности пациентов.

\subsection{Постановка гипотез}
На основании выявленных пробелов настоящее исследование формулирует следующие гипотезы.

\textbf{Гипотеза 1.}
\begin{itemize}
    \item H0: предпочтение формата консультации и уровень дохода независимы.
    \item H1: респонденты с высоким уровнем дохода значимо чаще отдают предпочтение исключительно живому врачу по сравнению с респондентами с низким доходом.
\end{itemize}

\textbf{Гипотеза 2.}
\begin{itemize}
    \item H0: стоимость приема не влияет на вероятность выбора формата «только ИИ».
    \item H1: при высокой стоимости приема респонденты значимо реже выбирают ИИ без участия врача, требуя за эти деньги участия человека.
\end{itemize}

\textbf{Гипотеза 3.}
\begin{itemize}
    \item H0: повышение уровня точности диагностики не увеличивает статистически значимо вероятность выбора альтернативы, или данный атрибут не имеет наибольший относительный вес среди всех атрибутов.
    \item H1: повышение уровня точности диагностики статистически значимо увеличивает вероятность выбора альтернативы, и данный атрибут имеет наибольший относительный вес среди всех атрибутов.
\end{itemize}

\textbf{Гипотеза 4.}
\begin{itemize}
    \item H0: опыт посещения дерматолога не влияет на вес атрибута точности при выборе формата консультации.
    \item H1: респонденты, обращавшиеся к дерматологу («Да, много раз» или «Да, несколько раз»), придают статистически значимо больший вес атрибуту точности по сравнению с респондентами, никогда не посещавшими дерматолога.
\end{itemize}

\textbf{Гипотеза 5.}
\begin{itemize}
    \item H0: респонденты, ответившие «4--5» на утверждение «Я доверяю врачу больше, чем алгоритмам ИИ», не демонстрируют статистически значимо более низкую вероятность выбора формата «только ИИ» или не демонстрируют более высокую полезность от формата «врач + ИИ».
    \item H1: респонденты, ответившие «4--5» на утверждение «Я доверяю врачу больше, чем алгоритмам ИИ», демонстрируют статистически значимо более низкую вероятность выбора формата «только ИИ» и более высокую полезность от формата «врач + ИИ».
\end{itemize}

\newpage

% Глава 2
\section{ДАННЫЕ И МЕТОДОЛОГИЯ}

\subsection{Структура опросника и дизайн DCE-эксперимента}
Для измерения потребительских предпочтений в отношении форматов дерматологической консультации с применением ИИ был разработан онлайн-опросник на основе метода дискретного выбора (Discrete Choice Experiment, DCE). Опрос распространялся через Google Forms методом снежного кома (snowball sampling) — через социальные сети и личные контакты участников исследования — в марте 2026 года. Итоговая выборка составила 296 респондентов.

Опросник включал пять содержательных блоков. Первый блок — социально-демографический: возраст, пол, уровень образования, тип населенного пункта, уровень владения цифровыми технологиями. Второй блок — финансовый: три вопроса о субъективной оценке материального положения, из которых в дальнейшем формировался составной индекс дохода для проверки гипотезы 1. Третий блок — медицинский: опыт обращения к дерматологу и наличие диагностированных кожных заболеваний. Четвертый блок — шкалы Ликерта об отношении к ИИ в медицине, в том числе ключевой для гипотезы 5 вопрос «Я доверяю врачу больше, чем алгоритмам ИИ». Завершали опрос 8 карточек DCE-выбора.

DCE-эксперимент был построен на четырех атрибутах дерматологической консультации, представленных в таблице 1.

\begin{table}[h]
\centering
\caption{Атрибуты и уровни DCE-эксперимента}
\label{tab:attributes}
\small
\begin{tabularx}{\textwidth}{l>{\raggedright\arraybackslash}X>{\raggedright\arraybackslash}X}
\toprule
\textbf{Атрибут} & \textbf{Уровни} & \textbf{Обозначение в модели} \\
\midrule
Тип провайдера & Только ИИ; Только врач; Врач + ИИ & provider (категориальная, база --- ИИ) \\
Точность диагностики & 80\%; 85\%; 90\%; 95\% & accuracy (непрерывная, доли) \\
Время ожидания & Немедленно; 1 сутки; 3 суток; 1 неделя & waiting\_time (0 / 1 / 3 / 7 дней) \\
Стоимость приема & 1 500 руб.; 3 500 руб. & price\_group (межгрупповой, не входит в MNL) \\
\bottomrule
\end{tabularx}
\end{table}

Стоимость приема варьировалась между группами, но не внутри одного набора карточек: все три альтернативы в каждом сете предъявлялись по одной и той же цене. Рандомизация ценовой группы осуществлялась через встроенный attention check (проверка внимания): респондентам предлагалось взглянуть на текущее время и указать, четная или нечетная сейчас минута. Попавшие в группу с четной минутой видели цену 1 500 руб., с нечетной --- 3 500 руб. Такой дизайн позволяет тестировать гипотезу 2 на уровне агрегированных групп, однако исключает включение цены в MNL-модель как варьирующего атрибута, поскольку все три альтернативы в каждом трайле имеют одинаковую стоимость и коэффициент цены в модели не идентифицируется.

Каждый респондент получал 8 карточек выбора (trials). В каждой карточке одновременно предъявлялись три альтернативы --- «Только ИИ», «Только врач» и «Врач + ИИ» --- с различными комбинациями уровней точности и времени ожидания. Атрибуты варьировались циклически в соответствии с заранее заданным ортогональным дизайном.

\subsection{Описание собранной выборки}
Выборка формировалась нецелевым методом снежного кома, что обусловило ее специфический демографический профиль. Медианный возраст респондентов составил 19 лет, средний --- 20,5 лет; диапазон --- от 16 до 67 лет. Подавляющее большинство участников (45,3\%) --- подростки 19 лет; доля респондентов старше 25 лет не превышает 5\%. Гендерное распределение: 62,5\% женщин и 37,5\% мужчин. По уровню образования 73,6\% имеют неоконченное высшее образование, что отражает преобладание студенческой аудитории. Географически 83,4\% респондентов проживают в Москве или Санкт-Петербурге.

Уровень цифровой грамотности выборки высокий: 46,3\% оценивают его как «выше среднего», 22,3\% --- как «продвинутый»; лишь единицы затруднились с использованием базовых приложений. По финансовому положению 61,5\% отметили, что денег на основные нужды вполне хватает; 34,5\% --- что хватает, но приходится экономить. Опыт обращения к дерматологу имеют 67,9\% респондентов: 50,7\% обращались несколько раз, 17,2\% --- много раз; 32,1\% никогда не посещали дерматолога. Наиболее распространенные диагностированные заболевания --- акне (30,4\%) и атопический дерматит / экзема / псориаз (22,3\%). У 38,2\% кожных заболеваний не диагностировано.

Указанные характеристики выборки необходимо принимать во внимание при интерпретации результатов: молодая городская аудитория с высокой цифровой грамотностью не является репрезентативной для российского населения в целом, что составляет одно из ключевых ограничений исследования.

\subsection{Очистка и контроль качества данных}
При предобработке сырые данные Google Forms прошли первичную фильтрацию и очистку в R с использованием пакета tidyverse. Из выборки были исключены респонденты, не прошедшие встроенную проверку внимательности (ожидаемый ответ --- 4 балла из 5), а также респонденты старше 27 лет и проживающие в населенных пунктах с численностью менее 100 тыс. жителей --- в соответствии с целевой аудиторией исследования.

Оставшиеся анкеты прошли проверку на внутреннюю согласованность ответов по системе штрафных баллов. Выделялось пять типов нарушений:
\begin{enumerate}
    \item противоречие между оценкой финансового положения и способностью покрыть непредвиденные расходы;
    \item одновременный выбор варианта «Никогда не диагностировались» и любого конкретного заболевания;
    \item наличие диагноза при отрицании каких-либо кожных проблем в настоящем и прошлом;
    \item одновременная максимальная оценка текущих и прошлых кожных проблем;
    \item большое расхождение между двумя вопросами о конфиденциальности данных.
\end{enumerate}
Анкеты с суммарным штрафным баллом 3 и выше исключались из анализа, так как содержали логически несовместимые ответы, свидетельствующие о недостаточной внимательности при заполнении.

Дополнительно выявлялись и удалялись случаи прямолинейного ответа (straightlining) --- когда респондент проставлял одинаковые оценки по всем шкалам Ликерта (стандартное отклонение < 0,5 или доля модального ответа > 80\%).

Итоговый датасет в длинном формате содержал 6~384 наблюдения (266 респондентов после очистки × 8 трайлов × 3 альтернативы).

\subsection{Методология анализа}
\subsubsection{Метод дискретного выбора: теоретические основы и процедура}
Метод дискретного выбора (Discrete Choice Experiment, DCE) представляет собой количественный инструмент, предназначенный для выявления и измерения предпочтений потребителей в ситуациях, где выбор между несколькими альтернативами определяется набором их характеристик (атрибутов). В отличие от традиционных опросных методов (например, шкал Лайкерта или ранжирования), DCE моделирует реальные компромиссы (trade-offs), с которыми сталкивается индивид при принятии решения. Респонденту последовательно предъявляются гипотетические сценарии (карточки выбора), в каждом из которых необходимо выбрать одну из двух или более альтернатив, различающихся по значениям атрибутов (например, точность диагностики, время ожидания, стоимость). Многократное повторение таких выборов позволяет эконометрически оценить относительную важность каждого атрибута и готовность индивида обменивать один атрибут на другой.

Теоретической основой DCE служит теория случайной полезности (Random Utility Theory, RUT), предложенная Д. Макфадденом (2000, Нобелевская премия). Согласно RUT, полезность, которую индивид $n$ получает от альтернативы $j$ в ситуации выбора $t$, может быть разложена на детерминированную (наблюдаемую исследователем) компоненту, зависящую от атрибутов альтернативы и, возможно, характеристик индивида, и случайную (ненаблюдаемую) компоненту, которая включает влияние неучтённых факторов, измерительные ошибки и индивидуальные вариации вкусов. Предположение о распределении случайной компоненты (чаще всего -- распределение Гумбеля (тип I экстремального значения)) приводит к мультиномиальной логит-модели (MNL), в которой вероятность выбора альтернативы $j$ имеет замкнутую форму:
\[
P_{n}(j) = \frac{\exp(V_{nj})}{\sum_{k=1}^{J} \exp(V_{nk})},
\]
где $V_{nj}$ --- детерминированная часть полезности, обычно задаваемая как линейная комбинация атрибутов. Коэффициенты этой линейной комбинации интерпретируются как вклад каждого атрибута в полезность (предельная полезность), а их соотношения позволяют рассчитать предельную норму замещения между атрибутами (например, готовность платить за повышение точности или сокращение времени ожидания).

Применение DCE в здравоохранительных исследованиях имеет ряд преимуществ. Во-первых, метод позволяет оценить спрос на гипотетические услуги или технологии, которые ещё не получили широкого распространения (как, например, автономная ИИ-диагностика). Во-вторых, DCE количественно измеряет вес каждого атрибута и выявляет неоднородность предпочтений, что невозможно при использовании простых опросов о намерении. В-третьих, результаты DCE могут быть непосредственно использованы для прогнозирования рыночных долей при изменении характеристик услуги (метод имитации выбора).

В настоящем исследовании DCE-эксперимент включал 8 карточек выбора на каждого респондента. В каждой карточке одновременно предъявлялись три альтернативы: «Только ИИ», «Только врач» и «Врач + ИИ». Атрибуты «Точность диагностики» (4 уровня: 80\%, 85\%, 90\%, 95\%) и «Время ожидания» (4 уровня: немедленно, 1 сутки, 3 суток, 1 неделя) варьировались независимо для каждой альтернативы в соответствии с ортогональным планом. Атрибут «Стоимость приема» был постоянным в пределах каждой карточки (межгрупповая рандомизация: 1500 руб. или 3500 руб.) и не включался в детерминированную компоненту полезности как варьирующий, но использовался для проверки ценовой чувствительности на агрегированном уровне (гипотеза 2). Такой дизайн позволяет избежать статистической неидентифицируемости коэффициента цены (поскольку в каждом выборе все три альтернативы имеют одинаковую цену) и одновременно протестировать, влияет ли уровень цены на частоту выбора формата «Только ИИ» между группами респондентов.

\subsubsection{Спецификации моделей}
В работе оценивались две основные спецификации.

\textbf{Спецификация 1 (базовая модель, гипотеза 3)} -- без альтернативно-специфических констант (ASC), с общими коэффициентами для атрибутов точности и времени ожидания и категориальным атрибутом провайдера, играющим роль ASC:
\[
V_{inj} = \beta_{acc} \cdot \text{accuracy}_{inj} + \beta_{wait} \cdot \text{waiting\_time}_{inj} + \delta_j,
\]
где:
\begin{itemize}
    \item $\text{accuracy}_{inj}$ -- точность диагностики для альтернативы $j$ (непрерывная, от 0,80 до 0,95);
    \item $\text{waiting\_time}_{inj}$ -- время ожидания в днях (дискретная: 0, 1, 3, 7);
    \item $\delta_j$ -- фиксированный эффект альтернативы $j$ (для $j =$ «Только врач» и $j =$ «Врач + ИИ»; для базовой альтернативы «Только ИИ» $\delta = 0$).
\end{itemize}
Коэффициенты $\beta_{acc}$ и $\beta_{wait}$ являются общими для всех трёх альтернатив, что соответствует дизайну, в котором точность и время ожидания могут варьироваться независимо для каждого провайдера.

\textbf{Спецификация 2 (модель с взаимодействием, гипотеза 5)} -- с альтернативно-специфическими константами и индивидуальной переменной доверия, взаимодействующей с этими константами:
\[
V_{inj} = \beta_{acc} \cdot \text{accuracy}_{inj} + \beta_{wait} \cdot \text{waiting\_time}_{inj} + \text{ASC}_j + \gamma_j \cdot \text{high\_trust}_n,
\]
где:
\begin{itemize}
    \item $\text{ASC}_j$ -- альтернативно-специфическая константа для $j =$ «Только ИИ» и $j =$ «Врач + ИИ» относительно базовой альтернативы «Только врач» (для которой $\text{ASC}=0$);
    \item $\text{high\_trust}_n$ -- бинарная переменная, равная 1, если респондент $n$ ответил «4» или «5» на утверждение «Я доверяю врачу больше, чем алгоритмам ИИ», и 0 в противном случае;
    \item $\gamma_j$ -- коэффициенты взаимодействия, показывающие, насколько изменится полезность альтернативы $j$ (относительно базовой) при переходе от группы с низким доверием к группе с высоким доверием.
\end{itemize}

Оценка параметров производилась методом максимального правдоподобия с использованием пакета \texttt{mlogit} языка R. Стандартные ошибки рассчитывались на основе гессиана. Качество подгонки оценивалось с помощью псевдо-$R^2$ Макфаддена. Для проверки гипотез о взаимодействии (гипотезы 4 и 5) использовались односторонние тесты (alternative = «greater» или «less» в соответствии с направлением гипотезы). Все расчёты выполнены на данных в длинном формате (6~384 наблюдения: 266 респондентов $\times$ 8 трайлов $\times$ 3 альтернативы).

\newpage

% Глава 3
\section{ЭМПИРИЧЕСКИЕ РЕЗУЛЬТАТЫ И ИХ ОБСУЖДЕНИЕ}

\subsection{Гипотеза 1: доход и предпочтение живого врача}
Для проверки гипотезы 1 был построен составной индекс дохода на основе трех финансовых вопросов опросника. Корреляции Пирсона между компонентами составили $r = 0,42$--$0,57$, коэффициент $\alpha$ Кронбаха --- 0,756, что свидетельствует о приемлемой внутренней согласованности шкалы. Разбивка по медиане индекса сформировала группы: High Income (96 респондентов) и Low Income (163 респондента). Ещё 7 человек были исключены из рассмотрения, так как выбрали вариант «Затрудняюсь ответить» на вопрос об общем финансовом положении. Медиана выбрана вместо среднего как более устойчивая к выбросам мера центральной тенденции. Из анализа с помощью $z$-теста были также исключены респонденты с мажоритарным выбором «Врач + ИИ», поскольку данный метод применялся к бинарной переменной. Баланс ценовых групп в разрезе групп дохода был проверен хи-квадратом: $p = 0,023$, перекос статистически значим, поэтому ценовая группа была включена в регрессию как контрольная переменная. Однако ее включение не изменило выводов: все $p$-value при переменной дохода остались существенно выше 0,05.

Доля выбора формата «Только врач» составила 97,75\% в группе высокого дохода и 98,65\% в группе низкого дохода. Результаты тестирования представлены в таблице 2.

\begin{table}[h]
\centering
\caption{Результаты проверки гипотезы 1}
\label{tab:hyp1}
\small
\begin{tabularx}{\textwidth}{p{3.5cm} X c p{2.2cm}}
\toprule
\textbf{Метод} & \textbf{Статистика} & \textbf{$p$-value} & \textbf{Вывод} \\
\midrule
Z-тест долей & $z = 0$ & 0,5 & H0 не отвергается \\
Мультиномиальная логит-регрессия \newline (коэф. дохода) & $\beta = -0,541$ (Доктор) (SE=1,022) \newline $\beta = -0,811$ (Доктор+ИИ) (SE=1,113) & 0,596 / 0,466 & H0 не отвергается \\
\bottomrule
\end{tabularx}
\end{table}

По всем примененным методам статистически значимых различий между группами не обнаружено. Коэффициенты регрессии при переменной дохода отрицательны для обеих альтернатив ($-0,541$ для «Только врач» и $-0,811$ для «Врач + ИИ» относительно базовой категории «Только ИИ»), что вопреки H1 указывает на то, что высокодоходные респонденты чуть реже выбирают живого врача. Однако $p$-value в обоих случаях существенно превышает 0,05, и данная тенденция статистически незначима.

Нулевая гипотеза не отвергается: уровень дохода не является значимым предиктором предпочтения живого врача. Содержательная интерпретация такова: в выборке формат «Только ИИ» выбирался крайне редко --- лишь 2 человека в каждой группе дохода. Когда почти все респонденты выбирают одно и то же, статистический тест не имеет достаточно вариации, чтобы обнаружить разницу между группами, даже если она существует. Это приводит к возможной ошибке второго рода: мы не отвергаем H0 не потому, что связи нет, а потому что данных для ее обнаружения недостаточно.

\subsection{Гипотеза 2: цена и выбор формата «Только ИИ»}
Для проверки гипотезы 2 сравнивались доли выборов формата «Только ИИ» в двух ценовых группах: при стоимости 1 500 руб. (группа «чётная минута», N = 134) и при 3 500 руб. (группа «нечётная минута», N = 125). Доля выборов «Только ИИ» составила 13,4\% в группе низкой цены (144 из 1 072 выборов) и 13,2\% в группе высокой цены (132 из 1 000 выборов).

$Z$-тест долей (односторонний, alternative = \texttt{"greater"}): $z = 0,156$, $p = 0,438$. Результат верифицирован смешанной логистической регрессией (GLMM) с контролем дохода и случайным эффектом на уровне респондента: коэффициент при ценовой группе составил $\beta = -0,0026$ (SE=0,184), $p = 0,989$ --- эффект статистически незначим и практически равен нулю.

Нулевая гипотеза не отвергается: стоимость приема не оказывает значимого влияния на вероятность выбора формата «Только ИИ». Это может объясняться тем, что принятие решения в пользу ИИ-диагностики определяется прежде всего установками доверия и воспринимаемой точностью, а не ценой.

\subsection{Гипотеза 3: точность как доминирующий атрибут}
Базовая MNL-модель оценивалась на полном датасете в длинном формате. Модель сошлась за 6 итераций; log-likelihood составил $-1\,377,3$. Для оценки качества подгонки был рассчитан Pseudo $R^2$ Макфаддена относительно модели только с константами: $R^2 = 0,324$. Это свидетельствует о хорошем качестве подгонки модели.

Результаты представлены в таблице 3.

\begin{table}[p]
\centering
\caption{Результаты MNL-модели и относительная важность атрибутов (гипотеза 3)}
\label{tab:hyp3}
\footnotesize
\setlength{\tabcolsep}{4pt}
\begin{tabular}{lcccc}
\toprule
\textbf{Атрибут} & \textbf{$\beta$} & \textbf{$p$-value} & \textbf{Диапазон полезности} & \textbf{Относительная важность} \\
\midrule
accuracy & 19,42 & $<2,2\times10^{-16}$ & 2,912 & 52,5\% \\
provider (doctor) & 1,796 & \multirow{2}{*}{$<2,2\times10^{-16}$} & \multirow{2}{*}{1,796} & \multirow{2}{*}{32,3\%} \\
provider (doctor\&ai) & 1,479 & & & \\
waiting\_time & $-0,121$ & $<2,2\times10^{-16}$ & 0,844 & 15,2\% \\
\bottomrule
\end{tabular}
\end{table}

Коэффициент при accuracy статистически значим на любом стандартном уровне значимости ($p < 2,2\times10^{-16}$) и положителен: повышение точности диагностики на 1 процентный пункт увеличивает логарифм шансов выбора данной альтернативы на 0,194. Коэффициент при waiting\_time отрицателен и значим: каждый дополнительный день ожидания снижает полезность альтернативы. Оба коэффициента при provider значимы и положительны относительно базовой категории «Только ИИ», при этом «Только врач» ($\beta = 1,796$) предпочтительнее «Врача + ИИ» ($\beta = 1,479$).

Расчет относительной важности показывает, что accuracy занимает первое место с долей 52,5\%, опережая provider (32,3\%) и waiting\_time (15,2\%). Таким образом, обе части гипотезы 3 подтверждаются: точность диагностики статистически значимо увеличивает вероятность выбора альтернативы и имеет наибольший относительный вес среди всех атрибутов. Этот результат согласуется с данными китайского DCE-исследования \cite{Liu2021}, где точность также занимала первое место (39,29\%).

\subsection{Гипотеза 4: опыт посещения дерматолога и вес точности}
Для проверки гипотезы 4 в MNL-модель был введен interaction term accuracy:visited, где visited = 1 для респондентов, обращавшихся к дерматологу хотя бы несколько раз (N = 183), и 0 --- для тех, кто никогда не посещал специалиста (N = 83).

Результаты модели: $\beta(\text{accuracy}) = 20,642$ (базовый эффект для группы visited = 0), $\beta(\text{accuracy:visited}) = -1,745$, стандартная ошибка взаимодействия = 1,415, $z = -1,233$, $p = 0,891$ (одностороннее).

Коэффициент взаимодействия отрицателен, то есть группа с опытом посещения дерматолога оценивает точность несколько ниже, чем не имеющая такого опыта --- что прямо противоположно H1. Однако эффект статистически незначим на любом стандартном уровне.

Нулевая гипотеза не отвергается: опыт обращения к дерматологу не влияет на вес атрибута точности при выборе формата консультации. Возможное объяснение состоит в том, что точность воспринимается как универсальный приоритет вне зависимости от клинического опыта пациента: обе группы потребителей одинаково ценят надежность диагноза как базовое требование к медицинской услуге.

\subsection{Гипотеза 5: доверие к врачу и предпочтение форматов}
Переменная high\_trust = 1 присвоена респондентам, ответившим «4» или «5» на утверждение «Я доверяю врачу больше, чем алгоритмам ИИ» --- всего таких в выборке 231 из 266 (173 с оценкой 5 и 58 с оценкой 4). MNL-модель с индивидуальной характеристикой high\_trust, взаимодействующей с альтернативными константами, оценивалась с базовой альтернативой «Только врач».

\begin{table}[h]
\centering
\caption{Результаты проверки гипотезы 5}
\label{tab:hyp5}
\small
\begin{tabular}{lcccc}
\toprule
\textbf{Коэффициент} & \textbf{$\beta$} & \textbf{SE} & \textbf{$z$} & \textbf{$p$-value (одностор.)} \\
\midrule
high\_trust : Только ИИ (H5.1) & $-1,125$ & 0,199 & $-5,66$ & $7,6 \times 10^{-9}$ \\
high\_trust : Врач + ИИ (H5.2) & $-0,519$ & 0,180 & $-2,876$ & 0,998 \\
\bottomrule
\end{tabular}
\end{table}

Гипотеза 5.1 подтверждается: коэффициент high\_trust при альтернативе «Только ИИ» равен $-1,125$ (SE = 0,199, $z = -5,66$, $p = 7,6\times10^{-9}$). Это означает, что респонденты с высоким доверием к врачу относительно ИИ имеют статистически значимо более низкую полезность формата «Только ИИ» при прочих равных атрибутах. Эффект является одним из наиболее робастных в исследовании и согласуется с глобальными данными: Wu et al. (2024) фиксируют, что 84,4\% пациентов предпочтут мнение врача мнению алгоритма в случае их расхождения.

Гипотеза 5.2 не подтверждается: коэффициент high\_trust при альтернативе «Врач + ИИ» равен $-0,519$ (SE = 0,180, $z = -2,876$, $p = 0,998$). Знак коэффициента противоположен ожидаемому: респонденты, высоко доверяющие врачу, оценивают полезность гибридного формата ниже, чем референтная группа. Этот результат можно интерпретировать следующим образом: для респондентов с сильным доверием к врачу гибридный формат «Врач + ИИ» воспринимается как излишнее усложнение --- им достаточно «Только врача», и добавление ИИ не увеличивает, а снижает субъективную ценность услуги.

\subsection{Общее обсуждение результатов}
Совокупность полученных результатов формирует последовательную картину потребительских предпочтений российской молодежной аудитории. Подавляющее большинство респондентов вне зависимости от дохода, цены и опыта обращения к врачу отдает предпочтение живому специалисту или гибридному формату --- формат «Только ИИ» выбирается лишь в 13\% случаев и не чувствителен ни к доходу, ни к цене. Этот результат воспроизводит глобальный консенсус, зафиксированный в научной литературе \cite{Jutzi2020, McRae2025}, и указывает на то, что барьер принятия ИИ носит не экономический, а психологический характер. Иными словами, снижение стоимости ИИ-консультации или рост доходов пациентов сами по себе не приведут к росту спроса на автономные ИИ-решения --- необходима работа с доверием и восприятием технологии.

Среди атрибутов услуги точность диагностики является бесспорным лидером по относительной важности (52,5\%), что согласуется с результатами зарубежных DCE-исследований \cite{Liu2021, Jagemann2024}. Это указывает на то, что потребители готовы принять ИИ-решения при условии гарантированно высокого качества диагностики --- то есть ключевым инструментом повышения спроса является не снижение цены, а развитие коммуникации с пациентами и повышение точности алгоритма.

Наконец, доверие к врачу является мощным предиктором отказа от формата «Только ИИ», однако не повышает привлекательность гибридного формата --- что свидетельствует о том, что для высокодоверяющих пациентов участие врача является необходимым и достаточным условием, а добавление ИИ воспринимается скорее как фактор неопределенности, чем как дополнительная ценность.

\newpage

% Заключение
\section{ЗАКЛЮЧЕНИЕ}
Настоящая работа была посвящена изучению потребительских предпочтений российской молодежной аудитории в отношении форматов дерматологической консультации с использованием искусственного интеллекта. Для достижения поставленной цели был разработан и проведен DCE-эксперимент на выборке из 296 респондентов, опрошенных в марте 2026 года, а результаты проанализированы с помощью мультиномиальных логит-моделей.

По итогам проверки пяти гипотез получены следующие результаты. Гипотезы 1 и 2 не подтверждены: ни уровень дохода, ни стоимость приема не оказывают статистически значимого влияния на выбор формата консультации. Гипотеза 3 подтверждена полностью: точность диагностики является наиболее значимым атрибутом с относительной важностью 52,5\%, статистически значимо увеличивая вероятность выбора альтернативы ($\beta = 19,42$, $p < 2,2\times10^{-16}$). Гипотеза 4 не подтверждена: опыт обращения к дерматологу не дифференцирует вес атрибута точности между группами. Гипотеза 5 подтверждена частично: высокое доверие к врачу значимо снижает вероятность выбора формата «Только ИИ» ($\beta = -1,125$, $p = 7,6\times10^{-9}$), однако не повышает полезность гибридного формата «Врач + ИИ».

Совокупность результатов формирует три содержательных вывода. Во-первых, барьер принятия автономного ИИ в дерматологии носит психологический, а не экономический характер: снижение цены или рост доходов пациентов не создает спрос на формат «Только ИИ». Главным драйвером принятия технологии является доверие, которое формируется через воспринимаемую точность алгоритма и сохранение участия врача в процессе. Во-вторых, точность диагностики является ключевым рычагом влияния на потребительский выбор --- именно ее повышение, а также улучшение коммуникации с пациентами, способны в наибольшей степени увеличить привлекательность ИИ-решений. В-третьих, гибридный формат «Врач + ИИ» не воспринимается как безусловно лучшая альтернатива для высокодоверяющих пациентов: для них участие живого врача является необходимым и достаточным условием, а добавление алгоритма не создает дополнительной ценности.

Полученные результаты согласуются с международными данными и вместе с тем вносят самостоятельный вклад: впервые применена методология DCE для изучения данной проблематики на российской выборке, что позволяет количественно сопоставить относительную важность атрибутов и выйти за рамки описательных опросных исследований.

Работа имеет ряд ограничений, которые следует учитывать при интерпретации результатов. Выборка сформирована методом снежного кома и смещена в сторону молодой городской аудитории с высокой цифровой грамотностью --- ее нельзя считать репрезентативной для российского населения в целом. Дизайн ценовой рандомизации через четность минуты порождает случайный, а не систематический дисбаланс групп, что снижает достоверность теста для гипотезы 2. Наконец, крайне низкая доля выборов формата «Только ИИ» создает риск ошибки второго рода при проверке гипотезы 1.

Перспективными направлениями дальнейших исследований являются:
\begin{itemize}
    \item репликация DCE-эксперимента на репрезентативной общероссийской выборке;
    \item включение ценового атрибута как варьирующего внутри трайла для полноценной оценки готовности платить;
    \item анализ гетерогенности предпочтений с помощью смешанных логит-моделей (Mixed Logit) с учетом социально-демографических ковариат;
    \item сравнительный анализ предпочтений в разрезе ценовых групп как отдельное направление, намеченное в настоящей работе.
\end{itemize}

% Список литературы
\begin{thebibliography}{99}
\bibitem{Traini2025} Traini et al., 2025. (Замените на полные данные)
\bibitem{Lamotkin2024} Lamotkin et al., 2024.
\bibitem{Jutzi2020} Jutzi et al., 2020.
\bibitem{McRae2025} McRae et al., 2025.
\bibitem{Artsrouni2025} Artsrouni et al., 2025.
\bibitem{Liu2021} Liu et al., 2021.
\bibitem{Jagemann2024} Jagemann et al., 2024.
\bibitem{Myzrova2024} Myzrova et al., 2024.
\bibitem{Shen2020} Shen et al., 2020.
\bibitem{Polesie2020} Polesie et al., 2020.
\bibitem{Orlova2023} Orlova et al., 2023.
\bibitem{VCIOM2024} ВЦИОМ, 2024.
\bibitem{Wu2024} Wu et al., 2024.
\bibitem{Kearney2025} Kearney et al., 2025.
\bibitem{Garcia2025} Garcia et al., 2025.
\bibitem{Davis1989} Davis, F.D. (1989). Perceived usefulness, perceived ease of use, and user acceptance of information technology. \textit{MIS Quarterly}, 13(3), 319–340.
\bibitem{Venkatesh2003} Venkatesh et al. (2003). User acceptance of information technology: Toward a unified view. \textit{MIS Quarterly}, 27(3), 425–478.
\end{thebibliography}

\end{document}